% ============================================================================
%  SECCIÓN 4.7: IMPLEMENTACIÓN DE ESTADO LOCAL
% ============================================================================

\section{Implementaci\'on de estado local}

El estado local se gestiona mediante el hook \textit{useState} de React,
que permite a los componentes mantener y actualizar su propio estado
interno. Cada vez que el estado cambia, el componente se re-renderiza
para reflejar la nueva situaci\'on en la interfaz.

\subsection{Estados en formularios}

Los formularios de inicio de sesi\'on y registro utilizan m\'ultiples
estados locales para gestionar los valores de los campos, los mensajes
de error y el estado de env\'io.

\subsubsection{Estado del formulario de registro}

El formulario de registro gestiona los siguientes estados locales:

\begin{lstlisting}
// Estado para errores del servidor
const [serverError, setServerError] = useState<string | null>(null);

// Estado para mensaje de exito
const [successMessage, setSuccessMessage] = useState<string | null>(null);

// Estado del formulario (react-hook-form)
const {
  control,
  handleSubmit,
  formState: { errors, isSubmitting },
} = useForm<RegisterFormData>({
  resolver: zodResolver(registerSchema),
  defaultValues: {
    fullName: '',
    email: '',
    password: '',
    confirmPassword: '',
  },
});
\end{lstlisting}

Cuando el usuario completa el formulario y lo env\'ia, la funci\'on
\textit{onSubmit} actualiza el estado seg\'un el resultado de la operaci\'on:

\begin{lstlisting}
const onSubmit = async (data: RegisterFormData) => {
  setServerError(null);
  setSuccessMessage(null);

  const { error } = await signUp(
    data.email, data.password, data.fullName
  );

  if (error) {
    setServerError(error);
    return;
  }

  setSuccessMessage(
    'Cuenta creada. Revisa tu email para confirmar tu cuenta.'
  );

  setTimeout(() => {
    router.replace('/auth/login');
  }, 2500);
};
\end{lstlisting}

\subsection{Estados en la pantalla de b\'usqueda}

La pantalla de cat\'alogo gestiona estados locales para los filtros y la
b\'usqueda, que se actualizan de forma independiente y reactiva:

\begin{lstlisting}
// Termino de busqueda (se actualiza inmediatamente)
const [searchQuery, setSearchQuery] = useState('');

// Termino debounced (se actualiza despues de 300ms)
const [debouncedQuery, setDebouncedQuery] = useState('');

// Filtro por tipo de escenario
const [selectedType, setSelectedType] = useState<string | null>(null);

// Filtro por deporte
const [selectedSport, setSelectedSport] = useState<string | null>(null);
\end{lstlisting}

Los filtros se combinan de forma acumulativa mediante \textit{useMemo},
lo que permite que m\'ultiples filtros se apliquen simult\'aneamente:

\begin{lstlisting}
const filteredScenarios = useMemo(() => {
  if (!scenarios) return [];

  return scenarios.filter((scenario) => {
    const matchesName = debouncedQuery
      ? scenario.nombre
          .toLowerCase()
          .includes(debouncedQuery.toLowerCase())
      : true;

    const matchesType = selectedType
      ? scenario.tipo?.toLowerCase() === selectedType.toLowerCase()
      : true;

    const matchesSport = selectedSport
      ? scenario.scenario_sports?.some(
          (ss) => ss.sports?.nombre?.toLowerCase()
            === selectedSport.toLowerCase()
        )
      : true;

    return matchesName && matchesType && matchesSport;
  });
}, [scenarios, debouncedQuery, selectedType, selectedSport]);
\end{lstlisting}

\subsection{Estado del bot\'on de favorito}

El bot\'on de favorito en la pantalla de detalle utiliza un estado local
para controlar la animaci\'on de carga mientras se ejecuta la operaci\'on
de toggle:

\begin{lstlisting}
const [imageError, setImageError] = useState(false);
\end{lstlisting}

Adem\'as, el hook \textit{useToggleFavorite} gestiona un estado interno
\textit{isToggling} que deshabilita el bot\'on durante la operaci\'on:

\begin{lstlisting}
const [isToggling, setIsToggling] = useState(false);

const toggleFavorite = async (userId: string, scenarioId: string) => {
  setIsToggling(true);
  try {
    // ... logica de favorito
  } finally {
    setIsToggling(false);
  }
};
\end{lstlisting}

\subsection{Ejemplo visual de cambio de estado}

Cuando el usuario presiona el bot\'on de favorito, el estado cambia y
la interfaz se actualiza visualmente:

\begin{center}
\begin{tikzpicture}[
  node distance=1.2cm,
  btn/.style={rectangle, draw=black!60, rounded corners, minimum width=3cm,
    minimum height=0.8cm, align=center, font=\small},
  flecha/.style={-Stealth, thick}
]
  \node[btn] (no-fav) {\textit{Guardar} (corazon vacio)};
  \node[btn, right=of no-fav] (toggling) {Cargando...};
  \node[btn, right=of toggling] (fav) {\textit{Guardado} (corazon lleno)};
  \draw[flecha] (no-fav) -- node[above] {onPress} (toggling);
  \draw[flecha] (toggling) -- node[above] {exito} (fav);
\end{tikzpicture}
\end{center}
